\section{Limitations}
\label{sec:limitations}

\subsection{Single 38-day calendar window, single exchange}
\label{sec:limitations-window}

This is the project's own most-flagged open item and is stated
explicitly rather than left implicit. Every result in this paper
describes BTCUSDT (and, for the transfer ladder, ETH/SOL/DOGE USDT)
perpetual-futures microstructure on Binance during a single 38-day
calendar window, 2024-02-22 -- 2024-03-30. No claim in this paper should
be read as asserting that market microstructure dynamics are
time-invariant, or that the specific magnitudes reported (the
$M_1$-vs-$M_2$ PR-AUC gap, the branching-ratio values, the transfer
ladder's captured-gap percentages) generalize to other calendar periods,
volatility regimes, or exchanges without further testing. This is a
standard scope limitation for microstructure econometrics, not a defect
specific to this study, but it bears directly on external validity: the
mechanisms studied here -- order-flow clustering, cross-excitation
between quote and trade events, queue depletion -- are not expected to
have fundamentally changed in kind since early 2024, but their exact
parameterization (kernel weights, branching ratios, the $M_1$-vs-$M_2$
PR-AUC gap magnitude) should not be assumed stable across different
market regimes (e.g.\ markedly higher or lower realized volatility, a
different funding-rate environment, or a different exchange's matching
engine and tick-size conventions) without a fresh fit. The
window itself was not a free choice: it is bounded on the recent end by
Binance's permanent discontinuation of the \texttt{bookTicker} feed for
BTCUSDT futures after 2024-03-30 (Section~\ref{sec:data-feeds}), so a
more recent replication under the identical protocol is not currently
possible against this exact data source. This was re-checked directly
rather than assumed stale: as of 2026-09-18 the live archive still lists
nothing past 2024-03-30, including a monthly rollup file one step past
the daily cutoff that turned out to contain only 6h47m of 2024-04-01
data. As the nearest feasible substitute, Section~\ref{sec:results-robustness}
reruns the core learning curve against an independent, non-overlapping
test week drawn from \emph{within} the same 38-day window and finds the
same $M_1>M_2$ result -- a within-window robustness check that
corroborates the finding's stability inside this window, but does not,
and cannot, speak to external validity across a genuinely different
calendar period or market regime.

\subsection{The excluded 2024-02-04 -- 2024-02-21 disorder window}
\label{sec:limitations-disorder}

A full-coverage audit of the originally planned 60-day window found
severe raw \texttt{bookTicker} ordering disorder between 2024-02-04 and
2024-02-21 -- on some days in this range, up to $\sim$40\% of a day's
rows showed a \texttt{transaction\_time} decrease relative to the
previous row (e.g.\ 18{,}728{,}015 of 46{,}396{,}761 rows on 2024-02-12),
a scale of disorder far beyond isolated same-millisecond ties that the
project's standard defensive sort (by \texttt{transaction\_time}, then
\texttt{update\_id}) produces \emph{a} valid-looking chronological order
for, but whose recovery of the \emph{true} underlying order at this scale
of scrambling was never independently verified. This 18-day span was
excluded from the primary window entirely rather than included with an
unverified repair, and no Hawkes fit, learning-curve point, or any other
result in this paper touches it. This is a conservative, disclosed
exclusion, not a silent gap: the resulting 38-day window is the largest
\emph{contiguous, independently re-audited clean} span available between
the resolved end of this disorder period (2024-02-22, confirmed via a
zero-raw-decrease re-audit) and the permanent \texttt{bookTicker}
discontinuation (2024-03-30). Whether the excluded period's underlying
market dynamics (which, unlike the archive's row ordering, were not
inherently anomalous) would change any of this paper's findings if
recovered and independently validated is untested and remains open.

\subsection{Cross-machine Hawkes non-bit-reproducibility}
\label{sec:limitations-reproducibility}

Several results in this paper (e.g.\ Table~\ref{tab:m3-comparison}'s
$M_2$ column differing from Table~\ref{tab:learning-curve} by
$\sim 0.0001$--$0.0002$ PR-AUC at every $N$, despite an identical
specification) reflect a finding that was itself investigated and
resolved during the project: the unconstrained L-BFGS-B Hawkes fit used
throughout this paper is \emph{not} bit-reproducible across different
machines or software environments, even at an identical fixed random
seed, identical data, and identical code -- most likely due to different
BLAS/LAPACK builds shifting the optimizer's floating-point trajectory.
This was distinguished directly from true non-determinism: two
back-to-back reruns of an identical 100{,}000-event calibration budget
\emph{on the same machine} were confirmed bit-identical to full displayed
precision, while the same specification run on a different machine in an
earlier session differed by up to $0.0016$ PR-AUC -- roughly an order of
magnitude larger than the previously assumed ``run-to-run noise floor''
of $\sim0.0001$--$0.0002$, which in retrospect was itself cross-machine
noise mislabeled as generic non-determinism. LightGBM fits ($M_0$/$M_1$,
which never call the Hawkes optimizer) reproduce exactly across every
run and every machine used in this project, at every reported precision.
Every cross-model or cross-configuration \emph{comparison} reported in
this paper (e.g.\ every $\Delta$ column, every GW statistic) uses model
outputs produced within a single run on a single machine, so this
non-reproducibility does not threaten any comparison's internal validity;
it does mean that a reader attempting to reproduce an individual $M_2$ or
$M_3$ PR-AUC figure to full displayed precision on different hardware
should expect a difference on the order of $10^{-4}$--$10^{-3}$, not
bit-exact agreement, and any future rerun on a \emph{single, fixed}
machine should reproduce exactly -- a discrepancy under those specific
conditions would indicate a real bug, not expected Hawkes noise.

\subsection{Transfer-ladder simplifications, resolved}
\label{sec:limitations-transfer}

The cross-asset transfer ladder (Section~\ref{sec:results-transfer}) was
disclosed, when first built, as resting on three simplifications relative
to the specification's stated protocol: transferring the full
$10\times10\times5$ excitation tensor rather than only the collapsed
branching-ratio matrix (a strict superset of what is needed, not a
simplification against the protocol); estimating the target asset's
baseline intensity $\boldsymbol\mu$ via a closed-form empirical rate
rather than a properly profiled maximum-likelihood re-estimation under
the fixed, transferred $\boldsymbol\alpha$; and transferring BTC's
$\boldsymbol\Gamma$ directly to each of ETH, SOL, and DOGE rather than
through a literal daisy chain (each asset's own re-estimated
$\boldsymbol\Gamma$ feeding the next rung). All three were subsequently
tested directly rather than left as caveats. The literal daisy chain was
built and evaluated and produced results within $0.0008$--$0.0009$ PR-AUC
of the direct-transfer design, with no consistent direction. A proper
fixed-$\boldsymbol\alpha$, profiled-$\boldsymbol\mu$ estimator (solved by
one-dimensional bisection, including an explicit boundary-clamp case
validated against the joint MLE's own boundary solutions to within
0.08\% relative error) was implemented and used to rerun every transfer
evaluation; every resulting PR-AUC changed by at most $0.0003$, with no
consistent direction across the eight cases tested. \textbf{These
simplifications are therefore resolved, not open}: the transfer-ladder
results reported in Section~\ref{sec:results-transfer} and
Table~\ref{tab:transfer-synthesis} reflect the originally disclosed
simplified design, and this subsection records that a substantially more
faithful re-implementation of the specified protocol was subsequently
confirmed to change nothing materially. What remains genuinely untested
for the transfer ladder is orthogonal to these three points: only one
BTC source calibration window was ever used as the transfer origin, and
no formal significance test (Giacomini--White or Romano--Wolf) was run on
any transfer rung -- every transfer number in this paper is a single
train/test point estimate on its target asset, unlike the BTC-only
results elsewhere in this paper, which are corroborated by a formal
rolling test. A further robustness extension to three more assets
(BNB, XRP, ADA, Section~\ref{sec:results-transfer}) also revealed a
new, unresolved question rather than only confirming the original
finding: those three rungs cluster $1.3$ percentage points apart at
$78.97\%$--$80.27\%$ gap-capture, below the original three-asset range's
floor of $81.7\%$, and the qty-tie dual-$\boldsymbol\Gamma$ check above
was not repeated for them. Whether this reflects a structural property
of the newly added assets or is a coincidence of six single-window,
single-source point estimates is not resolved in this paper.

\subsection{Other disclosed scope decisions}
\label{sec:limitations-other}

A small number of additional scope decisions were made explicitly during
the project and are recorded here for completeness rather than treated
as open questions requiring resolution in this paper. The post-hoc
stationarity projection (Section~\ref{sec:methodology-hawkes}) was
retained rather than replaced with an in-fit constraint, on the basis of
an empirical scan across 185 fits in which the projection essentially
never bound, later found to bind in exactly one hourly window that traces
to a known, pre-existing \texttt{bookTicker} data gap
(2024-03-08) rather than to a general large-$N$ property; the affected
window contributes 0.116\% of the $N=31$d training rows and a $0.47\%$
coefficient-level rescale for that one hour. A residual $8.4\%$
tie-order sensitivity in the fitted $\boldsymbol\Gamma$ matrix, after
masking the dominant same-millisecond touch-price/trade ambiguity
(Section~\ref{sec:data-tiebreak}), was traced to a narrower, disclosed
class of same-millisecond quantity-only quote/trade ties (6.35\% of
events) that the current masking scope does not cover; this scope
decision was made explicitly and re-reviewed (and reconfirmed adequate)
before both the $M_3$ build and the transfer ladder, each of which
consumes $\boldsymbol\Gamma$ in a different way. The difference-in-differences
placebo control specified for validating $M_3$'s structural claim
(Section~\ref{sec:methodology-placebo}) was not run in this study: given
$M_3$'s null PR-AUC result on the Primary Endpoint
(Section~\ref{sec:results-m3}), it was judged to have comparatively
little to falsify at the time relative to isolating the source of $M_3$'s
subsequently discovered calibration effect via ablation, which was run
instead and is reported in Section~\ref{sec:results-m3-ablation}; running
the placebo control against the calibration effect specifically (rather
than the null PR-AUC result it was originally scoped against) remains a
natural next step. Finally, the secondary/exploratory grid's
\texttt{target\_a} implementation does not take a depletion-fraction
argument, so the nominal 18 named specifications in
Table~\ref{tab:romano-wolf} contain only 12 distinct tests; this is a
codebase quirk disclosed at the time it was discovered, not silently
corrected, and all 18 named columns are reported as specified.
